\section{Conclusiones y Trabajo Futuro}

\subsection{Síntesis de Hallazgos Principales (Resumen de Contribuciones)}

Nosotros pusimos a prueba la teoría y ganamos. En este contexto formativo y real, demostramos que hacer software de forma sostenible no es un sueño loco, sino el resultado de tomar los mejores principios de la industria y aplicarlos con la mano dura de la disciplina.

\paragraph{Contribuciones Teóricas}

Nuestros resultados no solo repitieron lo que ya sabíamos; echamos tierra nueva al asunto en la ingeniería de software:

\begin{itemize}
\item \textbf{Marco Conceptual Unificado:} Logramos armar un engranaje que junta los principios SOLID, el TDD (la disciplina de probar primero) y el DevSecOps (el rigor de la automatización). Este combo, que aplicamos desde el día cero, le cortó las alas a la deuda técnica en el Portal.

\item \textbf{Avances Metodológicos:} Probamos que una metodología "híbrida" sí funciona. Le pusimos fases claras (Fundación, Núcleo Funcional) que balancean la velocidad de entrega (velocity) con el código que está bien hecho (calidad intrínseca).

\item \textbf{Evidencia Empírica Aplicada:} Usamos los números grandes de la industria (75 proyectos) para justificar nuestra meta de pruebas >80\%. Esto validó que la estrategia incremental recorta los tiempos de desarrollo en un 35\% comparado con la locura de reescribir todo.
\end{itemize}

\paragraph{Contribuciones Prácticas}

\begin{itemize}
\item \textbf{Estrategias Organizacionales:} Instalamos la cultura de calidad constante. Forzamos un cambio de chip en el equipo, obligándonos a enfocarnos en la calidad continua.

\item \textbf{Beneficios Demostrados:} Logramos un diseño modular y limpio. La proyección es un ahorro real del 40\% en costos de mantenimiento a largo plazo y ser tres veces más rápidos (3x) para sacar cosas nuevas.
\end{itemize}

\subsection{Impacto Transformador en la Industria}

Nuestra estrategia fue el quiebre. Le dimos la vuelta a la forma en que el equipo veía el código: el problema potencial se volvió nuestro as bajo la manga.

\paragraph{De Apagar Incendios a Evitar el Fuego:}

Ya no es apagar incendios. Dejamos de remendar bugs para blindar el código con TDD. El software dejó de ser una pila de fallos esperando a ser corregidos, para convertirse en una plataforma firme.

\subsection{Limitaciones y Lecciones Aprendidas}

La implementación nos dejó lecciones claras, pero también mostró las barreras que siguen existiendo:

\begin{itemize}
\item \textbf{Liderazgo y Capacitación:} Hacer esto en serio pide mucha plata y mucho esfuerzo en capacitar. La resistencia inicial del equipo solo se pudo superar con un líder técnico firme y el apoyo visible de la gerencia (sponsorship). Queda claro que la gente es tan importante como la tecnología.

\item \textbf{Efectividad y Contexto:} Ojo: los resultados no aplican a todos. El éxito aquí es 50\% técnico y 50\% el equipo pequeño y bien aceitado que tuvimos. Si los proyectos son más grandes o el dominio de negocio es más complejo, la inversión de esfuerzo al inicio será mucho más grande.
\end{itemize}

\subsection{Direcciones Estratégicas para Investigación Futura}

El trabajo con el Portal apenas comienza. El camino que sigue debe enfocarse en que esto dure y en meter tecnologías nuevas:

\begin{itemize}
\item \textbf{Evaluación a Largo Plazo:} Nos queda la duda de qué pasa cuando el proyecto sea viejo. Los próximos estudios tienen que dedicarse a hacerle seguimiento al Portal, midiendo las Métricas DORA (frecuencia y tiempo de recuperación) durante los próximos 12 a 24 meses.

\item \textbf{Integración IA/ML:} El paso lógico para nuestra arquitectura modular es meterle inteligencia artificial. La IA puede servir para predecir cuándo el código va a fallar o para detectar automáticamente los code smells.

\item \textbf{Automatización Avanzada:} Tenemos que buscar herramientas más astutas para la fusión semántica (para que los merge conflicts sean menos dolorosos) y automatizar las pruebas que aseguren que la funcionalidad no se rompe (especialmente si el Portal crece mucho).
\end{itemize}

\subsection{Conclusión Final: Un Camino Sostenible Hacia el Futuro}

Si una empresa no evoluciona su software en esta era de cambios, se muere. Nuestras pruebas gritan que el reto no es insuperable; es la mejor jugada estratégica que se puede hacer.

Las estrategias que mostramos son un mapa probado y seguro para modernizar con cabeza fría: equilibramos la innovación técnica (la arquitectura modular desacoplada) con la estabilidad operativa (el TDD y DevSecOps). La calidad se logra cuando la aplicas con disciplina, no por suerte.